\documentclass[12pt,a4paper]{scrartcl}

\usepackage[utf8]{inputenc}
\usepackage[ngerman]{babel}
\usepackage[babel,german=guillemets]{csquotes}
\usepackage[backend=bibtex8]{biblatex}
\usepackage{amsmath}
\usepackage{amsfonts}
\usepackage{amssymb}

\bibliography{../../literatur/quellen}

\author{Emanuel Seidinger}
\title{Besprechungsprotokoll}
\subject{$\alpha$-Shapes Applikation für mobile Plattformen}
\date{13.\,03.~2013}

\begin{document}
\maketitle

\section*{Teilnehmer}
\begin{list}{}{}
\item apl. Prof. Dr. Christian Icking
\item Dr. Lihong Ma
\item Emanuel Seidinger
\end{list}

\section*{Übersicht}
\begin{itemize}
\item Funktionalität
\item Bedienkonzept
\item Entwicklungsumgebung
\item Deployment
\item Zu klärende Punkte
\end{itemize}

\section*{Funktionalität}

\subsection*{Existierende Applikationen}
Zur Erstellung von $\alpha$-Shapes existieren bereits Applikationen in Form von Java Applets.
Diese erfordern die Installation eines Java Plugins im Browser.
Ein Beispiel eines solchen Applets findet man unter \cite{alpha_applet}.

Die Eingabe einer Punktmenge erfolgt per Klick auf die linke Maustaste. Mit einem Druck auf die
rechte Maustaste lassen sich Punkte wieder entfernen. Punkte lassen sich per Drag and Drop verschieben.
Der Parameter $\alpha$ ist hier als Kehrwert des Radius derjenigen Scheiben definiert, welche zur
Konstruktion der $\alpha$-Shape verwendet werden. Er kann mit Hilfe eines Schiebereglers im Intervall 
$(-\infty, \infty)$ variiert werden, um unterschiedliche $\alpha$-Shapes zu erzeugen.

Zusätzlich zur $\alpha$-Shape lassen sich die folgenden graphischen Elemente einblenden:
\begin{itemize}
\item Den Rand derjenigen verallgemeinerten Kreisscheiben mit Radius $\frac{1}{\alpha}$, auf denen die
$\alpha$-Nachbarn liegen
\item Das Voronoi Diagramm der Punkte mit kleinsten bzw. größten Abstand
\item Die Delaunay Triangulation der Punkte mit kleinsten bzw. größten Abstand
\end{itemize}

\subsection*{Die zur erstellende Applikation im Vergleich zu bestehenden Lösungen}
Zur Umsetzung soll auf Technologien zurückgegriffen werden, welche moderne Browser ohne zusätzliche
Plugins interpretieren können. Dieses Projekt soll mit HTML5 realisiert werden. Es soll nur HTML, CSS und
JavaScript verwendet werden.

Zur besseren Anschaulichkeit, soll der Parameter $\alpha$ der Radius derjenigen verallgemeinerten Kreisscheiben
sein, welche zur Konstruktion der $\alpha$-Shape verwendet werden.

Zusätzlich zur bestehenden Lösung, soll die neue Applikation dem Benutzer markante Werte für $\alpha$ aus
dem $\alpha$-Spektrum zur Auswahl anbieten. Diese markanten Punkte stellen Rasten auf dem Schieberegler
zur Einstellung von $\alpha$ dar.

\section*{Bedienkonzept}
Da die Applikation auf mobilen Geräten ausgeführt werden soll, ist darauf zu achten, dass eine sinnvolle
Touch-Gestensteuerung möglich ist. Die Benutzeroberfläche muss auch an den Viewport angepasst sein.

Als Hilfestellung zur Einstellung von $\alpha$ wird für positive Werte von $\alpha$ eine Kreisscheibe mit variablen
Radius dargestellt, welche so auf der Zeichenfläche verschoben werden kann, dass sie keinen Punkt der Punktmenge
verdeckt.
Für negative Werte von $\alpha$, wird ein Ring mit variablen Durchmesser dargestellt, welcher auf der Zeichenfläche
verschoben werden kann, dabei aber immer alle Punkte der Punktmenge enthält.

\section*{Entwicklungsumgebung}
Als Entwicklungsumgebung wird Netbeans verwendet. Das Projekt wird als HTML5 Projekt angelegt. Für die Durchführung
der Unit-Tests wird Jasmin \cite{jasmin_test} verwendet. Der Quellcode wird in einem Subversion Repository verwaltet.

\section*{Deployment}
Die Applikation soll auf verschiedenen mobilen Plattformen laufen. Dies lässt sich durch PhoneGap \cite{phone_gap}
realisieren.

Die Applikation kann als Web-Applikation über den Browser gestartet werden können. Soll die Applikation Zugriff
auf die Hardware der jeweiligen Ausführungsplattform erhalten, muss die Web-Applikation in eine native Applikation
eingebettet werden. Dadurch können unter anderem Punktmengen gespeichert und geladen werden.

Zur Einbettung in eine Desktop-Applikation könnte hierfür der JavaFX WebView \cite{javafx_webview} verwendet werden.
Eine weiter Möglichkeit ist die Webkit Integration in QT \cite{qt_webkit}.

\section*{Zu klärende Punkte}
\subsection*{Sprungstelle bei $\alpha = 0$}

In der Seminararbeit ist als $\alpha$-Hülle einer Punktmenge $S$ für negative $\alpha$, wobei der Betrag von
$\alpha$ kleiner als der Radius der kleinsten Umkreisscheibe, aber größer als $0$ ist, die gesamte Ebene angegeben.

Die $\alpha$-Hülle für negative $\alpha$ ist gegeben durch:

\begin{equation}
	AH(S) = \bigcap \lbrace D(C, -\alpha) \mid C \in \mathbb{R}^2 \wedge D \cap S = S \rbrace
\end{equation}

Für das oben angegeben Intervall von $\alpha$ ergibt dies folgende $\alpha$-Hülle:

\begin{equation}
	AH(S) = \bigcap \lbrace \rbrace = \oslash
\end{equation}

Dies entspricht der leeren Menge und nicht der gesamten Ebene, wie oben angegeben.

Entsprechend findet man im Artikel \cite{alpha1983} unter Abschnitt II. Basic Notions folgenden Textstelle:
\begin{quotation}
Thus we have a family of $\alpha$-Hulls for $\alpha$ ranging from $-\infty$ to $\infty$. Sample members of this
family are \textbf{the entire plane (for $\alpha$ sufficiently large)}, the smallest enclosing circle of $S$ (when 
$\dfrac{1}{\alpha}$ equals its radius), the convex hull of $S$ (for $\alpha = 0$), and $S$ itself (for
$\alpha$ sufficiently small).
\end{quotation}

Hier ist die Originaldefinition der $\alpha$-Hülle für positive $\alpha$ aus \cite{alpha1983}:
\begin{quotation}
Definition 1: Let $\alpha$ be a sufficiently small but otherwise arbitrary positive real. The $\alpha$-hull of $S$ is
the intersection of all closed discs with radius $\dfrac{1}{\alpha}$ that contain all the points of $S$.
\end{quotation}

Für große $\alpha$ gibt es keine Kreisscheibe mehr, welche alle Punkte aus $S$ enthält. Die $alpha$-Hülle ergibt sich
demnach als Durchschnittsmenge der leeren Menge, welche wiederum die leere Menge ist.

\printbibliography

\end{document}
